\documentclass[../00_Main.tex]{subfiles}
\begin{document}

Funding constraints regarding the manufacturing of the conductive cranial shell led to an adaptation of the validation strategy, which originally aimed for a full physical integration of the software with the anatomical phantom head. Consequently, the system evaluation was conducted using a modular verification framework. Instead of a single final assembly test, each subsystem was isolated and tested independently to ensure strict compliance with the functional requirements:
\begin{itemize}
    \item The \textbf{Software Subsystem} was evaluated for latency and waveform spectral purity.
    \item The \textbf{Hardware Transduction} was validated using gelatine cubes and controlled input signals.
    \item The \textbf{Signal Acquisition Subsystem} was tested for signal integrity and noise levels using controlled input signals.
    \item The \textbf{Closed-Loop System} was assessed by integrating the signal generation and a simplified form of the hardware module, comparing the output using the acquisition module to verify real-time performance.
\end{itemize}

The validation process utilized both laboratory instrumentation and the developed software subsystem to generate controlled input signals.

To validate the behavior of the complete integrated system in the absence of the final physical assembly with the anatomical phantom head, the project utilizes the \textbf{Mathematical Model} as the integration reference for a two-layer phantom head constructed following the methods described in Section~\ref{sec:materials_and_methods}. This theoretical validation demonstrates that, given the proven performance of the individual components, the complete system functions according to the design specifications when the modules are combined.

\subsection{Gelatine Experiments}\label{subsec:gelatine_experiments}

\nota{WIP}

\subsection{Signal Generation Experiments}\label{subsec:signal_generation_experiments}

\nota{WIP}

\subsection{Signal Acquisition Experiments}\label{subsec:signal_acquisition_experiments}

\nota{WIP}

\subsection{Closed-Loop Experiments}\label{subsec:closed_loop_experiments}

\nota{WIP}

\subsection{Mathematical Model Experiments}\label{subsec:math_model_experiments}

\dots


\biblio % Needed for referencing to working when compiling individual subfiles - Do not remove
\end{document}
